%% CS2640 Modern Storage Systems — Midterm Checkpoint
%% USENIX format, ~8 pages

\documentclass[letterpaper,twocolumn,10pt]{article}
\usepackage{usenix-2020-09}
\usepackage{graphicx}
\usepackage{booktabs}
\usepackage{amsmath}
\usepackage{listings}
\usepackage{multirow}
\usepackage{tabularx}
\usepackage[shortlabels]{enumitem}
\usepackage[normalem]{ulem}
\usepackage{float}

\lstset{
  basicstyle=\ttfamily\scriptsize,
  breaklines=true,
  frame=single,
  backgroundcolor=\color{gray!8},
  rulecolor=\color{gray!40},
  xleftmargin=4pt, xrightmargin=4pt,
}

\newcommand{\sysname}{AdaptiveCache}
\newcommand{\para}[1]{\smallskip\noindent\textbf{#1.}\ }
\newcommand{\etal}{\textit{et al.}}

\begin{document}
\date{}

\title{\Large \bf AdaptiveCache: Layout-Aware Context Management\\
       for Prefix-Cache-Efficient LLM Agents}

\author{
  {\rm Vlad Cainamisir}\\
  Harvard John A. Paulson School of Engineering and Applied Sciences\\
  {\small CS2640: Modern Storage Systems\quad---\quad Midterm Checkpoint, April 2026}
}

\maketitle

% =====================================================================
\begin{abstract}
LLM agents accumulate context across hundreds of tool calls. Compaction is
inevitable---but every existing approach (FIFO eviction, summarization, RL
compression) leaves survivors in a configuration that makes future requests
pay full recompute cost, because the prefix byte sequence has changed. We
argue that compaction should be a \emph{layout optimization problem}: choose
not just what to keep but where to place survivors, so that the prefix byte
sequence seen by the next request is identical to the last, and its KV states
are already cached. \sysname pins stable, high-value blocks at fixed prefix
positions and evicts volatile blocks from the suffix, amortizing one layout
reorganization across many future cache hits.

We report two experiments with honest discussion of their limitations. A
simple Anthropic-API prototype confirms that naive message-level eviction
(FIFO, compaction) costs more than no eviction at all---cache invalidation
outweighs token savings. A controlled vLLM experiment shows that suffix-first
eviction holds prefix cache hit rate at $\sim$99\% through eviction, while
prefix-first eviction (FIFO) collapses to $\sim$16\%, but suffers from warm
GPU cache contamination. We then describe what Anthropic's production system
(Claude Code) and the concurrent MEMENTO paper reveal about the space, and
conclude that the heuristic scoring approach is insufficient---the real
contribution should be a data-driven scorer trained on attention traces at
scale.
\end{abstract}

% =====================================================================
\section{Introduction}

Modern LLM inference servers implement \emph{prefix caching}: if consecutive
requests share a byte-identical prefix, KV states for those tokens are reused,
reducing effective compute from $O(N)$ to $O(1)$ for cached positions. For a
long-running agent, prefix caching can eliminate 80--95\% of prompt compute.
The requirement is strict: any modification to early tokens invalidates all
downstream cached KV states.

Context management is therefore a \emph{storage layout problem}. Every
agent context eventually exceeds its budget and must be compacted. The
question is not just \emph{what to keep}, but \emph{how to arrange survivors}
so that the resulting prefix byte sequence is identical to the previous
request, and its KV states are already cached in the server. Compact
correctly and you pay one layout reorganization amortized across many
future steps. Compact naively and every subsequent step pays full
recompute---making compaction self-defeating.

This is the \textbf{layout gap}: every existing system optimizes what to
keep but ignores where survivors land. KV eviction
methods~\cite{zhang2023h2o,xiao2024streamingllm,li2024snapkv,
liu2023scissorhands,cai2024pyramidkv} scatter survivors at arbitrary
positions---a different configuration per step, no cross-step prefix reuse.
Summarization systems~\cite{wu2025resum,lu2025supo,openhands2025}
generate new tokens, creating a byte sequence with zero cached prefix.
RL-trained approaches~\cite{zhou2025mem1,yu2025memagent,sun2025contextfolding}
require retraining and still replace the prefix at each turn.

\sysname closes the layout gap: identify stable, high-value blocks and pin
them at fixed prefix positions; evict volatile blocks from the suffix by
leaving holes, never renumbering; never summarize as a primary mechanism.
One layout reorganization establishes the right prefix; future steps append
to the suffix and get continuous cache hits until the importance structure
changes enough to warrant another reorganization.

% =====================================================================
\section{Related Work}
\label{sec:related}

\subsection{KV-Cache Eviction (Inference-Layer)}

These systems reduce inference memory by evicting KV entries during generation.
They share the layout gap: correct identification of important tokens, but no
cross-step prefix stability.

\para{H2O~\cite{zhang2023h2o}} Attention weights follow a power law: a small
``heavy hitter'' set captures most mass. H2O maintains this set plus a
recent window, achieving $1.9\times$ throughput. \sysname borrows the
cumulative attention signal.

\para{StreamingLLM~\cite{xiao2024streamingllm}} Initial tokens receive
disproportionate attention as a SoftMax dump (``attention sinks''). Keeping
4 sink tokens + a rolling window enables infinite streaming at $22.2\times$
speedup. This doubly motivates \sysname's stable prefix: initial tokens are
both semantically load-bearing and empirically critical.

\para{ScissorHands~\cite{liu2023scissorhands}} Proves that $>$95\% of
high-attention tokens in the first half of generation are also high-attention
in the second half (\emph{persistence of importance}). This justifies
\sysname's use of importance variance as a stability signal---consistently
important blocks can be predicted and pinned early.

SnapKV~\cite{li2024snapkv} and PyramidKV~\cite{cai2024pyramidkv} extend this
with per-head and per-layer selection, and Liu \etal~\cite{liu2023lostinmiddle}
show empirically that content at fixed early positions is attended to far
more reliably than middle-scattered content---a behavioral complement to the
caching argument. All remain per-query with no cross-step stability.

\subsection{Agent Context Management}

\para{Summarization-based} ReSuM~\cite{wu2025resum} is a plug-and-play
external summarizer (+4.5\% over ReAct, training-free). SUPO~\cite{lu2025supo}
trains a GRPO summarization policy (+14\% BrowseComp-Plus).
OpenHands~\cite{openhands2025} deploys nine composable condensation strategies
in production. All generate new summary tokens, invalidating the prefix cache.
ReSuM training-free is \sysname's primary comparison target.

\para{RL-trained state replacement} MEM1~\cite{zhou2025mem1} replaces the
entire context with a fresh internal state each turn via PPO ($3.5\times$
memory reduction, $3.7\times$ task performance). MemAgent~\cite{yu2025memagent}
maintains a 1024-token overwrite memory via Multi-Conv DAPO, extrapolating to
3.5M tokens. Both require training; applying either training-free consistently
hurts performance---a direct warning for \sysname's heuristic scoring approach.

\para{Context Folding~\cite{sun2025contextfolding}} Trains an RL agent to
branch into sub-trajectories and fold them upon completion (62\%
BrowseComp-Plus, 58\% SWE-bench). It operates at the sub-trajectory level,
not the token level, and is not prefix-cache-aware.

\subsection{Summary}

\begin{table*}[t]
\centering\small
\setlength{\tabcolsep}{8pt}
\begin{tabular}{lccc}
\toprule
\textbf{System} & \textbf{Training required?} & \textbf{Layout-aware?} & \textbf{Prefix-cache-aware?} \\
\midrule
FIFO / LRU                              & No       & No  & No \\
H2O / SnapKV / ScissorHands / PyramidKV & No       & No  & No \\
ReSuM~\cite{wu2025resum}                & Optional & No  & No \\
MEM1~\cite{zhou2025mem1}                & Yes      & No  & No \\
MemAgent~\cite{yu2025memagent}          & Yes      & No  & No \\
Context Folding~\cite{sun2025contextfolding} & Yes  & No  & Partial \\
MEMENTO~\cite{kontonis2026memento}      & Yes      & No  & No$^\dagger$ \\
\midrule
\textbf{\sysname (this work)}           & \textbf{No} & \textbf{Yes} & \textbf{Yes} \\
\bottomrule
\end{tabular}
\\[4pt]
\small{$^\dagger$ MEMENTO explicitly disables prefix caching; see \S\ref{sec:concurrent}.}
\caption{\sysname is the only system that is layout-aware and prefix-cache-aware without requiring training.}
\label{tab:comparison}
\end{table*}

% =====================================================================
\section{System Design}

\subsection{Block Segmentation and Scoring}

\sysname parses the agent's OpenAI-format message list into typed
\emph{blocks}: SYSTEM, TASK, THOUGHT, ACTION, and four observation types
(OBS\_FILE, OBS\_SHELL, OBS\_GREP, OBS\_ERROR). Each block is scored on
two independent axes:

\para{Importance} A weighted sum of: (a) \emph{structural type prior}---file
reads score 0.6, shell outputs 0.3, system/task 1.0; (b) \emph{reference
count}---how many times the block's identifiers (file paths, function names)
appear in subsequent thoughts; (c) \emph{cumulative attention}
(100\% milestone stub, requires real weights from the model).

\para{Stability} Whether importance is consistent over time: (a)
\emph{type stability prior}---file reads are stable (0.7), error messages
volatile (0.1); (b) \emph{importance variance}---low variance over recent
steps implies the block can safely be pinned.

The combined \emph{pin score} $p = \text{importance} \times \text{stability}$
drives zone assignment. The \emph{eviction priority}
$e = \text{importance} \times (1-\text{stability})$ evicts hot-but-volatile
blocks before cold-but-stable ones.

\subsection{Three-Zone Layout and Eviction}

Blocks are assigned to three zones: PIN ($p > \theta_\text{pin}$),
MIDDLE ($p > \theta_\text{mid}$), and SUFFIX (everything else). Eviction
fires SUFFIX-first when the total token count exceeds budget, cascading to
MIDDLE if needed, and never touching PIN.

The PIN zone is reconstructed in fixed order at the prefix on every step---
system prompt, task description, and high-value file reads at identical
absolute positions. The byte sequence at the prefix is therefore identical
across consecutive steps, enabling continuous prefix cache hits.

A layout reorganization fires every $\sim$15 steps (or when the top-$K$
blocks change composition by $>$30\%), invalidating the cache for one step
and establishing a new stable layout thereafter.

\subsection{Three-Tier Runtime Architecture}

The system operates three tiers at different timescales:

\begin{table}[h]
\centering\small
\setlength{\tabcolsep}{5pt}
\begin{tabular}{clcc}
\toprule
\textbf{Tier} & \textbf{Operation} & \textbf{Cost} & \textbf{Freq.} \\
\midrule
1 & KV block deletion  & $\approx$free & Every step \\
2 & Layout reorganization & Moderate & $\sim$15 steps \\
3 & Summarization (fallback) & Expensive & Emergency \\
\bottomrule
\end{tabular}
\caption{Three-tier pipeline. Goal: never reach Tier 3.}
\label{tab:tiers}
\end{table}

\para{Tier 1 --- every step} Low-scoring blocks are marked for eviction and
their KV entries deleted from LMCache via \texttt{POST /delete\_blocks}.
Because modern LLMs use post-rotated RoPE keys baked in at compute time,
evicted positions become unattendable without invalidating any downstream
KV state~\cite{zhang2023h2o,li2024snapkv}. No new tokens, no LLM calls,
no recomputation.

\para{Tier 2 --- every $\sim$15 steps} When the top-$K$ block composition
changes by $>$30\%, or accumulated holes exceed 40\% of allocated positions,
the layout optimizer re-sorts surviving blocks into PIN $\to$ MIDDLE $\to$
SUFFIX order. This invalidates the prefix for one step; the new layout is
then stable for the next $\sim$15 steps.

\para{Tier 3 --- emergency only} If eviction alone cannot meet the budget
(e.g., the PIN zone is overflowing), a full LLM summarization replaces
accumulated context. This always invalidates the prefix; the goal is to
never reach this tier.

\subsection{A Critical Note on the Heuristic Scorer}

The scoring framework described above is entirely heuristic. No signal in
the current implementation is grounded in empirical data: the type priors
(file reads score 0.6, shell outputs 0.3) are educated guesses, the
reference-counting signal relies on regex identifier matching, and the
stability signals are structural proxies with no validation. This scorer
was produced by the AI coding assistant and, on reflection, is not a
heuristic the author finds convincing.

The 0/5 SWE-bench patch submission rate confirms this directly: the scorer
evicts file reads the agent needs, forces re-reads, and increases cost.
The heuristics are not wrong in direction but far too coarse for the
decisions that actually matter---motivating the data-driven pivot in
\S\ref{sec:reflection}.

% =====================================================================
\section{Concurrent Work and Production Evidence}
\label{sec:concurrent}

\subsection{Claude Code Compaction~\cite{claudecode2025}}

Claude Code's compaction source code became publicly accessible in early
2026. It reveals a four-stage cheapest-first pipeline:
\textbf{Snip} (sliding window) $\to$
\textbf{Microcompact} (tool result clearing, no LLM call) $\to$
\textbf{Context Collapse} (sub-problem folding) $\to$
\textbf{Autocompact} (full LLM summarization, last resort).
This is the same tier hierarchy \sysname independently arrived at.

The most significant finding is \emph{cached microcompact}: when the
server-side prompt cache is warm, Claude Code enqueues \texttt{cache\_edits}
--- API instructions to surgically delete KV entries for specific old tool
results \emph{without invalidating the rest of the prefix}. This is
exactly \sysname's Tier 1, running at production scale:

\begin{lstlisting}
return {
  messages,  // local messages UNCHANGED
  compactionInfo: { pendingCacheEdits: {
    deletedToolIds: toolsToDelete, ... } }
}
\end{lstlisting}

Limitations relative to \sysname: targets only tool results; requires an
Anthropic-internal API; no layout step. \sysname generalizes all three.
Autocompact uses a \emph{forked agent} that inherits the parent's cached
prefix, so only the summarization request is sent---a clean solution to the
cost-of-compaction problem also not present in academic work.

\subsection{MEMENTO}

MEMENTO~\cite{kontonis2026memento} (Microsoft Research, 2026) trains models
to segment their own chain-of-thought into blocks, compress each into a
memento, and physically remove the thinking block's KV entries mid-generation
via a vLLM fork. It achieves 2--3$\times$ peak KV reduction and
$1.75\times$ throughput by enabling more concurrent requests in GPU memory.

MEMENTO and \sysname are orthogonal: MEMENTO compresses within-step
reasoning traces; \sysname manages the across-step agentic context. Critically,
MEMENTO \emph{disables} prefix caching---its memento KV states are enriched
with information from masked thinking blocks, violating the ``same tokens
$\Rightarrow$ same KV values'' assumption prefix caching requires.
\sysname's layout optimization is the complement: it preserves the byte
stability that MEMENTO deliberately sacrifices.

MEMENTO's key finding---that memento KV states implicitly encode information
from masked blocks (removing this channel drops AIME24 accuracy by 15
pp)---suggests a novel scoring signal for \sysname: blocks heavily attended
to by subsequent context have already propagated their information forward
and are cheaper to evict.

% =====================================================================
\section{Implementation}
\label{sec:impl}

\subsection{Two-Prototype Evolution}

The implementation went through two distinct phases. The first was a
lightweight Anthropic API prototype using \texttt{cache\_control}
prompt-caching breakpoints, built to validate the core hypothesis quickly:
does eviction hurt cache hit rate? Seven strategies were implemented and
benchmarked on SWE-bench, confirming the hypothesis and producing the
motivating results in \S\ref{sec:eval}. The second phase is the full
system, introducing a GPU inference server (vLLM + LMCache), a KV block
management API, and a three-tier eviction pipeline operating below the
message level.

\subsection{Core Library}

The \sysname library (\texttt{src/adaptive\_cache/}) is infrastructure-independent
and fully unit-tested (41 tests across 5 modules):

\begin{center}
\small\setlength{\tabcolsep}{4pt}
\begin{tabular}{ll}
\toprule
\textbf{Module} & \textbf{Responsibility} \\
\midrule
\texttt{segmenter.py} & OpenAI messages $\to$ typed \texttt{Block} objects \\
\texttt{scorer.py}    & 2D importance$\times$stability scoring \\
\texttt{evictor.py}   & Budget cascade: SUFFIX$\to$MIDDLE, never PIN \\
\texttt{layout.py}    & Zone assignment; prefix-stable reordering \\
\texttt{kv\_controller.py} & Three-tier controller; calls LMCache API \\
\bottomrule
\end{tabular}
\end{center}

Block types, zone assignments, and scorer signals are validated with synthetic
fixtures and hand-labeled examples. The library runs in CPU containers and
communicates with the GPU server via Modal method calls.

\subsection{GPU Inference Server}

The server runs on a single A100-80GB GPU via Modal. It spawns vLLM as an
OpenAI-compatible subprocess on port 8000 and
waits up to 10 minutes for the health endpoint to respond. LMCache integrates
as a \texttt{LMCacheConnectorV1} plugin, which intercepts the KV transfer
path. The internal management API is exposed on port 6999 via
\texttt{LMCACHE\_INTERNAL\_API\_SERVER\_ENABLED=True}, providing
\texttt{/delete\_blocks}, \texttt{/list\_blocks}, and \texttt{/reset}
endpoints.

Measuring real prefix cache hit rates required polling vLLM's Prometheus
\texttt{/metrics} endpoint and reading
\texttt{vllm:gpu\_prefix\_cache\_hit\_rate} after each request. The
\texttt{num\_cached\_tokens} field from the OpenAI-compatible API returned
\texttt{null} in vLLM 0.8.5 (it is only populated in 0.9+), which required
the indirect Prometheus approach.

\subsection{Compatibility Engineering}

Getting vLLM 0.8.5 and LMCache 0.4.3 to coexist required resolving three
distinct issues.

\para{LMCache build from source} The PyPI package for LMCache 0.4.3 did not
include CUDA kernels for the A100's compute capability (sm\_80). LMCache had
to be built from source with
\texttt{TORCH\_CUDA\_ARCH\_LIST='7.5;8.0;8.6'} and
\texttt{SETUPTOOLS\_SCM\_PRETEND\_VERSION=0.4.3} (to avoid git-based version
detection failure inside Modal's build environment). This added significant
image build time and required careful layer caching in the Modal image
definition.

\para{ABI mismatch (C++ shim)} LMCache 0.4.3's CUDA bindings called
\texttt{c10::cuda::c10\_cuda\_check\_implementation} with an \texttt{unsigned int}
line-number parameter. PyTorch 2.6.0 changed the signature to \texttt{int}.
The symbol was resolved at load time but the ABI mismatch caused a segfault
on the first CUDA error check. Fixed with a 10-line C++ shim compiled into a
shared library and \texttt{LD\_PRELOAD}'d at container startup:

\begin{lstlisting}[language=C++]
namespace c10 { namespace cuda {
  void c10_cuda_check_implementation(
    int e, const char* f,
    const char* n, unsigned int l, bool b) {
    c10_cuda_check_implementation(
      e, f, n, static_cast<int>(l), b); }}}
\end{lstlisting}

\para{vLLM 0.8.5 / LMCache 0.4.3 API mismatch} LMCache was written against
vLLM 0.9+, which passes an \texttt{engine\_id} string to the KV connector
at construction. vLLM 0.8.5 does not. This caused an assertion error at
startup. Separately, when the LMCache lookup client was \texttt{None}
(fallback mode), the vLLM scheduler called a method that assumed it was
populated, causing a second crash mid-request. Both were patched in LMCache
source before installation:

\begin{lstlisting}[language=Python]
# factory.py: supply default engine_id
metadata.engine_id = (
    metadata.engine_id or 'default-engine-0')

# vllm_v1_adapter.py: graceful degraded mode
if self.lookup_client is None:
    return 0, request_data
\end{lstlisting}

The patches are applied by \texttt{patch\_lmcache.py} as part of the Modal
image build, so the deployed container always has them without requiring a
forked repository.

\subsection{SWE-bench Harness}

The harness runs mini-swe-agent with a custom \texttt{AdaptiveCacheModel}
subclass that intercepts the message list before each API call. Four policies
are implemented (\texttt{none}, \texttt{fifo}, \texttt{adaptive},
\texttt{kv\_adaptive}). The harness runs 20 CPU containers in parallel (5 per
policy), each cloning the relevant SWE-bench repository at the correct base
commit and running the agent locally---avoiding Docker-in-Docker issues with
Modal.

SWE-bench instance metadata (300 instances) is bundled as a local JSON file
(\texttt{swebench\_instances.json}), eliminating a runtime HuggingFace API
dependency that caused flaky failures in early runs. Results are written to a
Modal volume and aggregated by a local entrypoint after all containers complete.

\subsection{Tier 1: The Unresolved Problem}

Tier 1 block deletion requires identifying which LMCache CPU store entries
correspond to evicted blocks. LMCache keys its store by a rolling SHA256 hash
chain over 256-token chunks:
\begin{align*}
h_0 &= \text{SHA256}(0_\text{le64} \,\|\, \text{chunk}_0)[:8] \\
h_i &= \text{SHA256}(h_{i-1} \,\|\, \text{chunk}_i)[:8]
\end{align*}
where token IDs are 4-byte little-endian. We re-implement this chain in
\texttt{compute\_block\_hashes()} and POST target hashes to
\texttt{/delete\_blocks} on port 6999.

Despite the infrastructure being wired end-to-end, the call consistently
returns \texttt{deleted=0}. The diagnostic tool \texttt{verify\_block\_hashes()}
lists all keys in LMCache's CPU store and checks overlap with our computed
hashes---but has not been successfully run in the deployed environment due to
Modal's subprocess isolation (LMCache's engine runs inside the vLLM worker
process, which is not directly accessible from the Modal method handler).
All cache hit rate results below come from vLLM's built-in GPU prefix cache,
not from LMCache block deletion.

The root cause of this isolation issue points toward a different
implementation strategy for the final prototype: rather than calling an
HTTP management API from outside the serving process, the right approach is
to modify LMCache and vLLM directly---either via a fork (as MEMENTO did for
its block masking) or by contributing the block-deletion interface upstream.
This would eliminate the process-boundary problem entirely and give direct
access to LMCache's internal state machine.

% =====================================================================
\section{Results}
\label{sec:eval}

\subsection{Scope and Methodology}

This project was primarily exploratory: build infrastructure, validate the
core mechanism, identify what a rigorous evaluation requires. The
implementation relied heavily on AI coding assistance (Claude Code), with the
trade-off that systematic issues were documented but not always blocked on.
A more carefully instrumented manual re-implementation is planned for the
final prototype.

The vLLM SWE-bench experiments were attempted but all trajectories show zero
steps completed---a container connectivity issue not investigated before the
deadline. SWE-bench task results below are from the Anthropic API prototype.
The vLLM stack produced one reliable result: the controlled cache hit rate
experiment in \S\ref{sec:vllm_exp}.

\subsection{Infrastructure Deployed}

The full vLLM serving infrastructure was built and operated successfully.
The harness dispatches up to 20 CPU agent containers in parallel (5 per
policy across 4 policies) against a single A100-80GB GPU server running
vLLM with LMCache. Each container clones the target repository locally,
runs the agent, and writes results to a shared volume. The GPU server
handled over 2,000 inference requests across all experiment runs, with
vLLM's built-in prefix cache achieving a 76.3\% server-side hit rate and
saving an average of 10,977 cached tokens per request. This confirms the
serving stack is functional at the required scale; the failure was in
orchestrating the end-to-end SWE-bench pipeline reliably from within
Modal containers, not in the inference server itself.

\subsection{Anthropic API Prototype}

The lightweight prototype used Anthropic's API with \texttt{cache\_control}
prompt-caching breakpoints. Seven strategies were compared on SWE-bench
instances with Claude Haiku 4.5. Each strategy ran as an independent agent
session, giving clean per-policy isolation.

\begin{table}[h]
\centering\small
\begin{tabular}{lrrr}
\toprule
\textbf{Strategy} & \textbf{Cost} & \textbf{Turns} & \textbf{Hit rate} \\
\midrule
full (no eviction) & \textbf{\$0.18} & \textbf{36} & 75.8\% \\
pairs              & \$0.26          & 40           & 67.4\% \\
fifo               & \$0.43          & 55           & 65.8\% \\
compaction         & \$0.62          & 51           & 49.6\% \\
\bottomrule
\end{tabular}
\caption{Anthropic API prototype, astropy-12907, Claude Haiku 4.5.
No-eviction is cheapest: eviction adds turns and cache misses that
amplify cost rather than reduce it.}
\label{tab:prototype}
\end{table}

The result is counterintuitive: keeping full context is consistently cheapest
and requires the fewest agent turns. FIFO costs $2.4\times$ more despite
evicting 30\% of the context. Message-level eviction changes the byte sequence
at the prefix, invalidating Anthropic's prompt cache---the savings on tokens
are outweighed by the cost of recomputing previously cached content. Compaction
is worst: the summary LLM call is expensive, the new summary tokens have no
cached prefix, and the agent needs more turns to recover from lost context.

Over 5 instances, the pattern holds: a summarization-based strategy
(episodes, not shown) achieves 85.5\% cache hit rate but costs
\$0.31/instance vs full context at \$0.23/instance, because it generates
more total tokens via summarization calls.

\subsection{Controlled vLLM Experiment}
\label{sec:vllm_exp}

We ran a controlled 10-step coding conversation on the vLLM server
(Qwen2.5-7B, budget=500 tokens, eviction forced at step 5).

\begin{table}[h]
\centering\small
\begin{tabular}{lcc}
\toprule
\textbf{Policy} & \textbf{Mean hit rate} & \textbf{Post-eviction} \\
\midrule
none (full context) & 94.4\% & --- \\
kv\_adaptive        & 98.6\% & 99.2\% \\
fifo                & 40.1\% & \textbf{15.7\%} \\
\bottomrule
\end{tabular}
\caption{vLLM prefix cache hit rate. Policies run sequentially; LMCache CPU
store reset between runs, but GPU cache is not (see text).}
\label{tab:hitrate}
\end{table}

\para{Limitation: warm GPU cache} Policies ran sequentially on the same GPU
server. \texttt{none} ran cold (step 1: 0\% hit rate, 158s latency).
\texttt{fifo} and \texttt{kv\_adaptive} ran with vLLM's GPU cache already
warm (step 1: 92.3\% hit rate, $<$1.2s latency). This inflates their mean
hit rates, making the mean comparison (98.6\% vs 94.4\%) unreliable.

\para{What the data does show} The post-eviction divergence is not explained
by warm cache. At step 5, FIFO drops to 15.7\%: it removes the oldest
messages, changing the prefix byte sequence---even warm cached tokens no
longer match. kv\_adaptive holds at 99.2\%: it removes the \emph{newest}
messages, keeping the oldest content at the prefix unchanged. This confirms
the core claim: evict from the suffix and the prefix cache survives; evict
from the prefix and it collapses.

\para{What is not tested} Tier 1 KV block deletion is broken (returns 0).
All hits are from vLLM's GPU prefix cache, not LMCache. The experiment
validates message-level suffix preservation, not the hole-leaving KV
mechanism.

\subsection{SWE-bench Task Results}

The SWE-bench task experiments using the vLLM server produced zero completed
trajectories due to the container connectivity issue described above. The
following results are from the Anthropic API prototype (Claude Haiku 4.5,
5 astropy instances, 8K budget):

\begin{table}[h]
\centering\small
\begin{tabular}{lrrrr}
\toprule
\textbf{Policy} & \textbf{Submitted} & \textbf{Cost} & \textbf{Steps} & \textbf{Hit rate} \\
\midrule
none (unlimited) & 5/5 & \$0.31 & 43 & 95.3\% \\
fifo (8K)        & 5/5 & \$0.40 & 51 & 53.6\% \\
adaptive (8K)    & 0/5 & \$0.17 & 30 & 61.1\% \\
\bottomrule
\end{tabular}
\caption{SWE-bench task experiment, Anthropic API (Haiku 4.5), 5 astropy
instances, 8K budget. Resolve rates unknown (\texttt{sb-cli} not run).}
\label{tab:swebench}
\end{table}

Adaptive submits no patches despite being cheapest (\$0.17) and fastest
(30 steps). The heuristic scorer evicts file reads the agent needs later,
causing re-reads, extra turns, and eventual failure to produce a patch.
FIFO costs 28\% more than full context but at least retains the content
the agent needs. These results confirm the scoring problem identified in
\S\ref{sec:reflection}: the mechanism is right, the scoring is not.

% =====================================================================
\section{Reflection and Future Directions}
\label{sec:reflection}

\subsection{The Scoring Problem}

The controlled experiment validates the \emph{mechanism}: suffix-first
eviction preserves prefix cache hits. The SWE-bench result (0/5 patches from
adaptive vs 5/5 from full context) exposes the \emph{scoring problem}: the
heuristic scorer does not know what to keep.

The current scorer uses structural type priors and reference counting---it
cannot distinguish a file read the agent will need in 10 steps from one it
glanced at once. Without this, the system evicts needed content, forces
re-reads, and increases cost. MEM1~\cite{zhou2025mem1} and
MemAgent~\cite{yu2025memagent} show the same failure mode: learned
compression applied training-free consistently hurts. Importance estimation
is empirical, not structural.

\para{The near-term pivot} Collect 100+ SWE-bench agent traces with
\texttt{--attention-backend EAGER}; derive oracle labels (did this block get
attended to later? did evicting it correlate with failure?); train a
lightweight scorer (MLP, pairwise ranking loss) on those labels; evaluate
against the heuristic baseline at equal task-success. The current
infrastructure is the right substrate. The heuristic system is the baseline
to beat.

\subsection{A Longer-Term Synthesis: Tetris + Folding + MEMENTO}
\label{sec:tetris}

Working through the related systems surfaced an insight that may be more
valuable than the current design. We describe it as a future research
direction.

\para{The Tetris observation} When an agent works on a coding task, it
progresses through a series of sub-problems: locate the bug, understand the
affected files, make an edit attempt, run the tests, handle the failure.
Each sub-problem generates a burst of context (error messages, file reads,
bash outputs, edit attempts). When the sub-problem is \emph{resolved}, that
entire burst becomes dead weight. If those blocks were placed in the suffix,
they can be evicted as holes at zero cost, leaving the prefix untouched.

The implication: if the rate of sub-problem completion roughly keeps pace
with new work starting, the context size is bounded at
$|\text{stable prefix}| + |\text{active sub-problem}|$, indefinitely. The
agent runs on a self-cleaning sliding window. Tier~3 summarization may never
fire.

\para{The limitation of existing approaches} Context Folding~\cite{sun2025contextfolding} already tracks sub-trajectory completion and folds finished
work. But folding generates summary tokens---a new byte sequence---that
invalidates the prefix cache at every fold event.
MEMENTO~\cite{kontonis2026memento} demonstrates in-place KV block eviction
at scale, but explicitly disables prefix caching because its KV states are
enriched with masked content.

Neither system achieves: sub-trajectory-aware eviction \emph{and} prefix
cache preservation simultaneously.

\para{The proposed synthesis}
\begin{enumerate}[1.,leftmargin=1.8em,topsep=2pt,itemsep=1pt]
  \item \textbf{Sub-problem tracking (from Folding):} identify sub-problem
    boundaries in the agent's trajectory---either via learned signals (like
    FoldGRPO) or structural heuristics (a git commit, a passing test, a
    completed tool chain). Place the active sub-problem's context in the
    suffix.
  \item \textbf{In-place KV eviction without summary tokens (from MEMENTO
    infrastructure):} when a sub-problem completes, use MEMENTO's vLLM fork
    to physically remove its KV blocks. No summary tokens are generated. The
    work product is in the repository; the conversation history of getting
    there is not needed.
  \item \textbf{Prefix stability (from \sysname):} the foundational context
    (task, key files, cross-sub-problem references) stays at fixed prefix
    positions throughout. Because we evict only suffix blocks and generate
    no new tokens, the prefix byte sequence is never touched---prefix caching
    remains active, unlike MEMENTO.
\end{enumerate}

\begin{figure*}[t]
\centering
\small
\setlength{\tabcolsep}{5pt}
\begin{tabular}{lll}
\toprule
\textbf{Phase} & \textbf{PIN prefix} & \textbf{SUFFIX} \\
\midrule
Sub-problem A active (steps 1--10)
  & \texttt{[sys][task][ast.py][utils.py]}
  & \texttt{[err\textsubscript{A}][bash\textsubscript{A}][fix\textsubscript{A}]\ldots} \\
A resolves: in-place KV eviction, no summary
  & \texttt{[sys][task][ast.py][utils.py]} \quad\textit{(unchanged)}
  & \texttt{[\sout{err\textsubscript{A}}][\sout{bash\textsubscript{A}}][\sout{fix\textsubscript{A}}]} $\leftarrow$ holes \\
Sub-problem B active (steps 11--20)
  & \texttt{[sys][task][ast.py][utils.py]} \quad\textit{(unchanged)}
  & \texttt{[err\textsubscript{B}][bash\textsubscript{B}][fix\textsubscript{B}]\ldots} \\
\bottomrule
\end{tabular}
\caption{Proposed Tetris synthesis across two sub-problems. PIN prefix bytes are
identical throughout all three phases---continuous KV cache hits. No summary tokens
are generated at any point. Context size stays bounded at
$|\text{prefix}| + |\text{current sub-problem}|$.}
\label{fig:tetris}
\end{figure*}

\para{What makes this different from both parents}
Context Folding pays a cache miss at every fold; \sysname avoids it by
evicting without generating summary tokens.
MEMENTO preserves implicit KV information from masked blocks; in the agentic
setting this is unnecessary---we are evicting genuinely completed work, not
reasoning we might still need.
The synthesis gets the sub-trajectory awareness of Folding, the in-place
KV eviction infrastructure of MEMENTO, and the prefix cache stability of
\sysname---a combination none of the three achieves individually.

\para{Open questions} Does the agent need to reference resolved
sub-problems? (Probably rarely for focused coding tasks; more often for
open-ended research tasks.) Can sub-problem boundaries be detected without
training, or does this require a FoldGRPO-style RL signal? How does this
compose with the learned scorer described above?

% =====================================================================
\section{Next Steps and Milestones}

\para{Immediate} Fix Tier 1 (verify hash scheme, validate hole-leaving demo);
run \texttt{sb-cli} to get resolve rates on existing patch files; sweep
budgets \{4K, 8K, 16K, 32K\} to find the Pareto frontier.

\para{Medium-term (100\% milestone)} Collect attention traces at scale
(\texttt{--attention-backend EAGER}, 100+ instances); derive oracle importance
labels; train and evaluate a learned scorer against the heuristic baseline.

\para{125\% milestone} GRPO-based reinforcement signal over eviction
decisions, or MEMENTO-style agent self-management for agentic (tool-use)
contexts.

\begin{table}[h]
\centering\small
\begin{tabular}{cp{5.5cm}}
\toprule
\textbf{Milestone} & \textbf{Status} \\
\midrule
75\%  & \textit{Substantially complete.} Core library (41 tests), baselines,
  harness, controlled experiment validating suffix eviction $\Rightarrow$
  prefix cache stability. \\[3pt]
100\% & \textit{Redefined.} Not more heuristic signals---attention trace
  collection + learned scorer trained on oracle labels. \\[3pt]
125\% & \textit{Redefined.} GRPO or MEMENTO-style agent self-management
  for tool-use contexts. \\
\bottomrule
\end{tabular}
\end{table}

% =====================================================================
\section{Conclusion}

The layout gap is real, and the mechanism for closing it works: suffix-first
eviction preserves prefix cache hits; prefix-first eviction destroys them.
What doesn't work is the heuristic scorer---confirmed by the SWE-bench
experiment, by the Anthropic API prototype, and by the analogous failure mode
in MEM1 and MemAgent. The Claude Code source confirms that surgical KV
deletion without prefix invalidation is production-ready. MEMENTO confirms
in-place masking at scale. The research direction pivots: from tuning heuristic
importance weights to collecting attention traces and training a learned scorer
grounded in what the model actually attends to. The current infrastructure is
the apparatus for that experiment.

\bibliographystyle{plain}
\bibliography{refs}

\end{document}
